% !TeX root = ../papers/nc-p3.tex
\section{六次残余の三重曲線と平面三次束の四つの変更}
\label{sec:sextic-residual-triple-curves}
\begin{theorem}\label{thm:sextic-residual-triple-curves}
重複平面 $H$ を二回除いた整六次残余 $G$ の三重点軌跡は有限である。
\end{theorem}
\begin{proof}[平面二次の収縮]
$H$ の外の三重印の集合を $P$ とする。$|P|\ge13$ であり $P$ を含む非零二次式はない
（その式と $H$ の式の積が二十印を消す）。一般平面切断の整六次曲線は算術種数十であり、
三重以上の各点で種数を少なくとも三使う。よって三重曲線の被約総次数は三以下である。

まず滑らかな平面二次曲線 $C$ が三重曲線とする。平面の式を $l$ とし、
平面座標 $u=(u_0,u_1,u_2)$ における $C$ の式を $q(u)$ とする。
$G\in(l,q)^3$ である。実際 $(l,q)$ は素な完全交叉で、冪の逐次商は完全交叉環上自由であり
冪には埋込素点がない。従って $(lu_0,lu_1,lu_2,l^2,q)$ の五二次式における三次式として書ける。
これは $(l,q)^3$ の各生成元の次数六の倍を単項式に分けても確認できる。
$q^3$ の係数は非零で、$l=0$ 上の $G$ の点は全て $C$ 上にある。
五二次式の像の閉包は $\mathbb P^4$ の滑らかな二次曲面体
$Q:\ z_3 z_4=q(z_0,z_1,z_2)$ で、$l\ne0$ では $z_3\ne0$ への同型である。
$G$ の対応する三次切断 $Y$ は整である。この開集合では既に整で、境界の二次錐面全体は
$z_4^3$ の係数が非零だから含まれない。

$Q$ 上の三重印二つを通る直線が $Q$ に含まれないとする。座標でこれらを $x,y$ 軸とし、
$Q$ の式を $xy+q_0(w)=0$ とできる。$w$ は三座標で $q_0$ は非退化である。
三次式の $xy$ を消し、第一軸の近くで $x=1,y=-q_0(w)$ とすれば、
三重条件により $x^3,x^2w,xw^2$ の項が全て消える。第二軸でも同様で、$Y$ の式は $F_3(w)$ のみとなる。
平面三次 $F_3$ は整。両軸の外で $x$ または $y$ が非零なら $[w]$ への射は滑らかで、
三次曲線は重複度二以下のため三重点はない。$x=y=0$ では $q_0(w)=0$ で、
その曲線の局所法座標が $-xy$ となる。$F_3$ が $[w]$ で滑らかなら、
接方向の微分が非零か、法方向が単元係数で現れ二次項 $xy$ が残るかである。
従って三重になる $[w]$ は整三次の特異点一つまでである。
ゆえにこの二印を含めても三印以下だが、$C$ 上は高々六印なので $l\ne0$ に七印以上を要する。
そこで全ての二印を結ぶ直線が $Q$ に入る場合だけが残る。極化によりこれらの点は全等方空間を張り、
$Q$ 上の一本の直線に載る。$[z_0:z_1:z_2:z_3]$ への線型射影で入力を復元するので、
七つの入力印も共線になって矛盾する。
\end{proof}
\begin{proof}[一本の三重直線に沿う三次式の族]
一本でも三重直線 $J$ があるとし、その零化二座標を $u=(u_0,u_1)$、残る二座標を $v$ とする。
$[u]$ の各方向の平面において入力を $(t u,v)$ と表示すれば、
$G(tu,v)/t^3$ は $(v_0,v_1,t)$ の三次式である。$J$ 外の印は十以上で、
各々は対応する平面三次の三重点である。この三次式はどの平面でも非零である。
同じ平面に相異なる三重点が二つあればその三次式は両点を結ぶ直線の三乗である
（二方向の微分が恒等零となる）。従ってどの群にも三印までで、異なる方向は四つ以上ある。

四つの異なる方向の印から非共面な四点を選べる。以下に線型代数の理由を記す。
$s$ 群から異なる群の独立な四代表が存在する条件は、任意の群集合 $I$ に対し
$r(I)+s-|I|\ge4$ である。ここで $r(I)$ はそれらの点のベクトル包の次元。
実際各群に同じ $s-4$ 次元の補空間の基底を加えると、全群の独立代表の条件に帰着する。
全群の場合の条件 $r(I)\ge |I|$ は帰納法で示せる。非空な真部分集合で等号なら先にそこで選び、
その包で割って残りに適用する。等号がなければ一群から非零元を選んで割っても残りの条件が保たれる。
元の条件のうち $I$ が全体の場合、全点が一平面に入れば $P$ はその平面と $J$ を含む平面との和に入るので不可能。
残りが一群なら $I$ が共線であっても合計六印までしかなく不可能。残りが二群なら $I$ に二方向以上あり
$r(I)\ge2$、残りが三群なら $I$ は非空、四群以上なら自明である。

方向の直線 $\mathbb P^1$ 上で $E=\mathcal O(1)^2\oplus\mathcal O(2)$ を取る。
各繊維の三成分を $v_0,v_1,t$ の係数とする。大域切断は $J$ を通る二次式であり、
上の六次式は $\operatorname{Sym}^3 E$ の切断である。実際 $t^j$ の係数の $u$ 次数は $3+j$ である。
選んだ四点での評価をそれぞれ一つの剰余体へ送り、その核束を $E'$ とする。
$\deg E'=0$ で、$H^0(E'(-1))=0$。$E(-1)$ の四切断は元の四座標の一次式に対応し、
四つの印は非共面だからである。直線上の束の分解より $E'$ は自明である。
四方向の一つで局所パラメータを $s$ とし、その印を消す二つの繊維一次式を $a,b$、残りを $c$ と取れば、
局所基底は $a,b,sc$ となる。全曲面での三重条件は $a^i b^j c^{3-i-j}$ の係数が
$s^{3-i-j}$ で割れることである（$c=1$ の局所座標を使う）。
ゆえに六次式は $\operatorname{Sym}^3 E'$ の切断であり、大域基底の三つの二次式 $q_i$ について
一つの定数係数整三次式 $F_3(q_1,q_2,q_3)$ に等しい。
四方向以外では $q_i$ が繊維の全一次座標であるから、その平面の $J$ 外には三重点がない。

各指定方向でも $J$ 外にさらに一印までしかないことを確かめる。
三つの $q_i$ による写像は局所的に、可逆な出力基底変更 $B(s)$ を用いて
$B(s)[a:b:sc]$ となる。中心以外の印は $s=0$ での像直線に行く。
この直線は $F_3$ の成分でない。別の三重印があれば平面上でも三重だから、その像は
$F_3$ の直線切断の唯一の三重根である。この根の逆像は中心を通る一本の直線である。
中心を除き $a$ または $b$ を単元にして正規化すると、$s$ 方向の出力直線に対する法微分は
$c$ と $a,b$ の一次式の和の非零定数倍である（$B(0)$ を除いて $B'(0)$ の寄与を加える）。
従ってこの逆像で微分階数が二未満になる点は一つだけである。階数二の点では平面への射が滑らかで、
整三次の重複度は二以下である。これで指定四平面でも各二印まで、合計八印までとなり十以上に反する。
\end{proof}
\begin{proof}[残る三次成分]
整平面三次を三重に含めば、その平面上で六次を超えるため平面因子が生じる。
非平面の整三次は有理正規三次 $C$ である。その三つの標準二次式は $C$ のイデアルを生成し、
$C$ の吹上げから二次式の値への射 $f$ は $\mathbb P^2$ 上の射影直線束である。
実際三つの二次式は二行三列Hankel行列の核座標であり、任意の非零核座標を指定した二式は
独立な一次式で一本の割線を与える。この入射はグラフの閉包で、二次式イデアルの吹上げである。
従って例外因子を $E_C$ とし入力超平面を $D$ とすれば $2D-E_C=f^*\mathcal O(1)$。
六次式は $3(2D-E_C)$ の切断なので平面の整三次の引戻しである。
$C$ 外は滑らかな $f$ の開部分であり、そこに三重点はない。
しかし $C$ 上の印は三次評価の全射性により十以下である。総次数三以下の全分割は以上で尽きた。
\end{proof}
